% The bundle-chain decomposition of a bacterial genome
% Build: pdflatex ecoli_decomposition (x3)
\documentclass[11pt]{article}
\usepackage{lmodern}
\usepackage[T1]{fontenc}
\usepackage{amsmath,amssymb,amsthm,mathtools}
\usepackage[margin=1.1in]{geometry}
\usepackage{booktabs}
\usepackage{microtype}
\usepackage[hidelinks]{hyperref}

\theoremstyle{plain}
\newtheorem{theorem}{Theorem}
\newtheorem{proposition}[theorem]{Proposition}
\newtheorem{lemma}[theorem]{Lemma}
\newtheorem{corollary}[theorem]{Corollary}
\theoremstyle{definition}
\newtheorem{definition}[theorem]{Definition}
\theoremstyle{remark}
\newtheorem{remark}[theorem]{Remark}

\newcommand{\Z}{\mathbb{Z}}
\newcommand{\K}{\mathcal{K}}
\newcommand{\coker}{\operatorname{coker}}
\newcommand{\Sk}{\mathrm{Sk}}

\title{\bfseries The bundle-chain decomposition of a bacterial genome:\\
\emph{Escherichia coli} K-12 at five $k$-mer lengths}
\author{Ruqing Chen\\[2pt]
\normalsize GUT Geoservice Inc., Rivi\`ere-Beaudette, Qu\'ebec J0P 1R0, Canada\\
\normalsize\texttt{ruqing@hotmail.com}}
\date{19 September 2026}

\begin{document}
\maketitle

\begin{abstract}
\noindent
The companion notes computed the sandpile group of a de Bruijn graph on a $5386$\,bp
phage. This one computes it on a $4\,641\,652$\,bp bacterium, \emph{Escherichia coli}
str.\ K-12 substr.\ MG1655 (NC\_000913.3), at $k=31,51,75,100,150$. The computation is
possible because of one observation that costs a line: in the convention used throughout
this series every vertex of $G_k$ satisfies $d^{+}(v)=d^{-}(v)=\operatorname{mult}(v)$, so
a vertex survives unitig contraction \emph{if and only if its $k$-mer is repeated}. The
$4.6$-million-vertex graph is therefore never built; the compacted graph is read off the
repeated $k$-mers directly, and the whole pipeline runs in about $13$ seconds. Three
findings. \emph{One:} the irreducible part of the ambiguity of a $4.6$\,Mbp genome is a
digraph on $115$ vertices at $k=150$ --- the chain $4\,641\,652\to23\,933\to115$ --- and
its group is computed exactly. \emph{Two:} the splitting is extremely lopsided; at
$k=150$ the length-driven summand has order $10^{10439}$ and the arrangement summand
$10^{49.7}$, so $99.5\%$ of the exponent is removable by longer $k$-mers and $0.5\%$ is
removable by no $k$ at all, while the irreducible part still grows by $10^{344}$ as $k$
falls to $31$. \emph{Three:} the seven ribosomal operons do \emph{not} form a
multiplicity-seven repeat family, at any $k$. Five lie on one strand and two on the
other, and the single-strand theory of this series sees a repeat only between co-oriented
copies; the two minus-strand operons give the longest clean spine in the genome
($4030$\,bp at $k=150$) and the five plus-strand ones a separate, fragmented family. This
is a quantitative measure of what a bidirected theory would have to supply. Six checks,
all exact, including a rerun of the whole fast pipeline against the published phage
decomposition.
\end{abstract}

\section{Introduction}\label{sec:intro}

For a circular sequence $S$ and $k\ge1$, $G_k(S)$ is the de Bruijn multigraph whose
vertices are the distinct $k$-mers and whose arcs are the occurrences of $(k{+}1)$-mers;
$S$ is an Eulerian circuit, and by the BEST theorem the number of reconstructions of $S$
from its $k$-mer spectrum is $|\K(G_k)|\prod_v(d^{+}(v)-1)!$ with $\K(G_k)$ the sandpile
group \cite{Bio1}. The bundle-chain splitting theorem \cite{Bio3} writes
\[
\K\bigl(G_k(S)\bigr)\;\cong\;\bigoplus_{\text{chains }c}(\Z/r_c)^{\,\ell_c-1}
\;\oplus\;\K(\Sk),
\]
a \emph{length-driven} summand that longer $k$-mers remove and an \emph{arrangement}
summand, the sandpile group of the repeat-adjacency skeleton, that no $k$ removes.
Everything so far has been computed on $\varphi$X174, $5386$\,bp. The question this note
answers is what the decomposition looks like on a genome three orders of magnitude
larger.

\paragraph{What is known.} Assemblers have compacted de Bruijn graphs to unitigs for
twenty years \cite{PTW,KM,BCALM}, and the repeat graph is the central object of the
long-read assembler Flye \cite{PTT,Flye}; the repeat content of \emph{E.\ coli} K-12 ---
seven \emph{rrn} operons, some fifty IS elements in a dozen families, several hundred
REP/BIME palindromic elements --- has been catalogued since the genome was published
\cite{Blattner}. What is new here is not the repeat catalogue but the algebra attached to
it: the exact invariant factors of the arrangement group, and the exact split between the
part of the ambiguity that a longer read length removes and the part it cannot.

\section{Why a 4.6\,Mbp genome is tractable}\label{sec:scaling}

\begin{lemma}[Branch vertices are the repeated $k$-mers]\label{lem:branch}
In $G_k(S)$, $d^{+}(v)=d^{-}(v)=\operatorname{mult}(v)$, the number of occurrences of the
$k$-mer $v$ in the circular sequence. Consequently a vertex has in- and out-degree one
--- i.e.\ is smoothed away by the unitig contraction of \cite[Lem.~4]{Bio1} --- if and
only if its $k$-mer occurs exactly once.
\end{lemma}

\begin{proof}
Each position $i$ contributes exactly one arc, from $S[i..i{+}k)$ to $S[i{+}1..i{+}k{+}1)$.
The arcs leaving $v$ are indexed by the positions at which $v$ occurs, and the arcs
entering $v$ by the positions at which $v$ occurs shifted by one; both sets have
$\operatorname{mult}(v)$ elements.
\end{proof}

This is immediate, and it is the whole algorithm. The compacted graph is obtained without
constructing $G_k$: let $R=(p_1<\cdots<p_m)$ be the positions at which a repeated $k$-mer
starts. Then the compacted multidigraph has one vertex per distinct repeated $k$-mer and
one arc $S[p_t..p_t{+}k)\to S[p_{t+1}..p_{t+1}{+}k)$ for each $t$ cyclically, because the
walk leaving a branch vertex runs through unique $k$-mers until the next repeat position.
For \emph{E.\ coli} this replaces $4\,641\,652$ vertices by between $23\,933$ and
$34\,272$ (Table~\ref{tab:main}), and the decomposition of \cite{Bio3} then runs in one
pass, since the chain-step relation is a partial injection \cite[Lem.~6]{Bio3}.

Finding the repeated $k$-mers is the only step that touches the whole genome. We hash
each of the $N$ positions to $64$ bits, take multiplicities with a sort, and then --- this
is what keeps the result exact --- recover the actual $k$-mer \emph{strings} at the
positions whose hash repeats and group by string. A hash collision can therefore only
add a candidate that is discarded; it cannot manufacture a repeat. No collision occurred
at any $k$ (check~(H)).

\section{Results}\label{sec:results}

\begin{table}[ht]
\centering\footnotesize
\begin{tabular}{rrrrrrrr}
\toprule
$k$ & distinct $k$-mers & branch & arcs & chains & skeleton
& $\log_{10}|\K(\Sk)|$ & $\log_{10}|\K(G_k)|$\\
\midrule
150 & 4\,594\,348 & 23\,933 & 71\,237 & 107 & \textbf{115} & 49.71 & 10\,488.89\\
100 & 4\,588\,709 & 26\,631 & 79\,574 & 166 & 175 & 71.11 & 11\,690.57\\
 75 & 4\,584\,636 & 28\,579 & 85\,595 & 243 & 261 & 99.80 & 12\,563.08\\
 51 & 4\,579\,164 & 31\,011 & 93\,499 & 342 & 366 & 143.17 & 13\,716.83\\
 31 & 4\,570\,869 & 34\,272 & 105\,055 & 683 & 870 & 393.75 & 15\,350.54\\
\bottomrule
\end{tabular}
\caption{\emph{E.\ coli} K-12 MG1655 under unitig contraction followed by bundle-chain
contraction. ``branch'' is the number of repeated $k$-mers (Lemma~\ref{lem:branch}) and
equals the vertex count of the compacted graph. All structural entries exact;
$\log_{10}|\K(\Sk)|$ by float LU, and exact as a group at $k=150,100$
(Table~\ref{tab:groups}).}
\label{tab:main}
\end{table}

\begin{table}[ht]
\centering\footnotesize
\begin{tabular}{rl}
\toprule
$k$ & $\K(\Sk)$, the arrangement group\\
\midrule
150 & $(\Z/2)^{17}\oplus(\Z/3)^{9}\oplus(\Z/4)^{12}\oplus(\Z/5)^{10}\oplus(\Z/7)^{3}
\oplus(\Z/8)^{4}\oplus(\Z/9)^{3}\oplus(\Z/16)^{3}$\\
    & $\qquad\oplus\,\Z/27\oplus(\Z/32)^{2}\oplus\Z/49\oplus\Z/21308303$\\[2pt]
100 & $(\Z/2)^{29}\oplus(\Z/3)^{16}\oplus(\Z/4)^{15}\oplus(\Z/5)^{15}\oplus(\Z/7)^{2}
\oplus(\Z/8)^{4}\oplus(\Z/9)^{2}\oplus(\Z/16)^{6}$\\
    & $\qquad\oplus\,\Z/25\oplus\Z/27\oplus(\Z/32)^{2}\oplus\Z/49\oplus\Z/81
\oplus\Z/128\oplus\Z/512\oplus\Z/3448649$\\
\bottomrule
\end{tabular}
\caption{The arrangement group, exactly, where the skeleton is small enough for the
$p$-adic Smith form (check~(K)). At $k=75,51,31$ the skeleton has $261$, $366$ and $870$
vertices and only the order is given; a $100$- to $400$-digit determinant could not be
factored in any case.}
\label{tab:groups}
\end{table}

\begin{remark}[The shape of the answer]\label{rem:shape}
Three things are worth reading off Table~\ref{tab:main}.

\emph{The skeleton is small.} At $k=150$ the entire arrangement ambiguity of a
$4.6$\,Mbp genome is carried by a digraph on $115$ vertices: $4\,641\,652\to23\,933$ by
unitig contraction, $23\,933\to115$ by bundle-chain contraction. The second step is the
one the theory supplies, and it is the larger reduction in proportion.

\emph{The split is lopsided.} At $k=150$ the length-driven summand has order
$10^{10439.18}$ and the arrangement summand $10^{49.71}$: $99.53\%$ of the exponent of
$|\K(G_{150})|$ is removable by reading longer $k$-mers, and $0.47\%$ is not removable by
any $k$. The percentage is not the interesting part --- $10^{49.7}$ is still an enormous
number of globally distinct threadings, and it is the part that only long reads, mate
pairs or a reference can resolve.

\emph{The irreducible part grows fast as $k$ falls.} From $k=150$ to $k=31$ the
arrangement summand grows from $10^{49.71}$ to $10^{393.75}$, a factor of $10^{344}$,
while the skeleton grows from $115$ to $870$ vertices. Shortening the read length does
not merely leave more length-driven ambiguity; it creates genuinely new arrangement
ambiguity, because repeats shorter than the old $k$ enter the graph and connect parts of
the skeleton that were previously unrelated.
\end{remark}

By the BEST theorem the number of circular sequences with exactly the $k$-mer spectrum of
this genome is $|\K(G_k)|\prod_v(d^{+}(v)-1)!$, which is $10^{20735.14}$ at $k=150$ and
$10^{31748.76}$ at $k=31$.

\section{The ribosomal operons, and a correction}\label{sec:rrna}

\emph{E.\ coli} K-12 has seven ribosomal RNA operons, each about $5$\,kb, and the natural
expectation --- the one the task that prompted this note recorded --- is a repeat family
of multiplicity seven. \emph{No such family exists, at any $k$.}

The reason is structural, not numerical. The NCBI feature table places five operons on
the $+$ strand (at $223\,771$, $3\,941\,808$, $4\,035\,531$, $4\,166\,659$, $4\,208\,147$)
and two on the $-$ strand (at $2\,726\,069$ and $3\,423\,423$). The theory of
\cite{Bio1,Bio2,Bio3,Bio4} is single-strand: two copies are a repeat of $G_k(S)$ only if
they are co-oriented \emph{in $S$}. The seven operons therefore cannot form one family.
They split by strand, and the computation confirms the split exactly (check~(R)):

\begin{itemize}
\item the \textbf{longest bundle chain in the whole genome}, at every $k$ tested, is the
multiplicity-\emph{two} spine of the two minus-strand operons --- $s=3881$ at $k=150$
(span $4030$\,bp), $s=4412$ at $k=51$ (span $4462$\,bp) --- annotated 16S, 23S and 5S
ribosomal RNA. Two copies annotated on the same (minus) strand are co-oriented with each
other, hence a direct repeat in $S$;
\item the five plus-strand operons give the multiplicity-\emph{five} family, whose longest
chain is the 16S gene (span $775$\,bp), with copies at $223\,997$, $3\,942\,034$,
$4\,035\,757$, $4\,166\,885$ and $4\,208\,373$ --- all five operons;
\item the five are not identical to one another either, so the rRNA region additionally
throws off chains of multiplicity three and four, one per subset of copies that agree.
\end{itemize}

Multiplicity seven does occur in this genome, but it is IS5A, not rRNA (five chains,
$88$ $k$-mers at $k=31$), as are multiplicities eight and nine. The other long chains are
all insertion sequences: IS186A ($\times3$, $1345$\,bp), IS2D ($\times3$, $1331$\,bp),
IS3B ($\times2$, $1260$\,bp), IS3A ($\times3$, $1255$\,bp), IS5U ($\times3$, $1202$\,bp).
At $k=31$ the highest multiplicities are $22$ and $23$ with spans of $32$ to $35$\,bp:
the REP/BIME extragenic palindromic elements, which are shorter than $51$ and so vanish
completely from the graph at every larger $k$ tested.

\begin{remark}[What this measures]\label{rem:bidirected}
A bidirected theory --- the open problem of \cite[Open (a)]{Bio2} --- would see all seven
operons as one family, and would see the IS copies that sit in inverted orientation. The
$5+2$ split is a quantitative statement of what the single-strand restriction costs on a
real genome: not a footnote, but the dominant feature of the largest repeat family
present. We do not have the bidirected theory and claim nothing about what it would give;
we can now say precisely what it would have to explain.
\end{remark}

\section{Verification}\label{sec:battery}

The script \texttt{ecoli\_sandpile.py} runs the six checks of Table~\ref{tab:battery} in
about $13$ seconds under half a gigabyte; its output is archived as
\texttt{ecoli\_sandpile\_output.txt}. The genome and the feature table are cached beside
it and re-fetched from the NCBI E-utilities API if absent.

\begin{table}[ht]
\centering\small
\begin{tabular}{clc}
\toprule
key & check & mode\\
\midrule
(V) & the fast path reproduces the published $\varphi$X174 decomposition, $k=9,\dots,12$ & exact\\
(H) & hashing finds candidates, strings decide; no collision survives & exact\\
(C) & the compacted graph is balanced with $d^{+}(v)=\operatorname{mult}(v)$ & exact\\
(B) & maximal bundle chains in one pass; Lemma~6 injectivity asserted on $1541$ chains & exact\\
(K) & the arrangement group, exactly for skeletons of $\le200$ vertices & exact/numeric\\
(R) & no rRNA chain has multiplicity $7$; the longest chain is the minus-strand pair & exact\\
\bottomrule
\end{tabular}
\caption{The verification battery, $6/6$.}
\label{tab:battery}
\end{table}

Check~(V) deserves a word. The pipeline used here is not the one used in \cite{Bio3}: it
never builds $G_k$, it contracts chains in a single pass instead of iteratively, and it
takes the Smith form $p$-adically instead of by Euclidean elimination. Run on
$\varphi$X174 it returns, at $k=9,10,11,12$, the same branch counts ($2$, $10$, $47$,
$146$), chain counts ($1$, $2$, $8$, $28$), skeleton sizes ($1$, $8$, $37$, $110$),
length-driven exponents and the same exact skeleton groups --- including
$\Z/4\oplus\Z/1076338579$ at $k=10$ and the seventeen elementary divisors at $k=9$. Three
independent implementations of the same theorem agreeing on a genome is the only
protection available when the target genome is too large to check by any other route.

\section{Scope}\label{sec:scope}

\emph{Exact.} Every structural quantity --- repeated $k$-mers, branch counts, the multiset
of $(r_c,\ell_c)$, skeleton sizes, genomic coordinates, the length-driven summand --- is
exact integer work, as is $\K(\Sk)$ at $k=150$ and $k=100$ and the location of every repeat
family against the feature table.

\emph{Numeric.} $\log_{10}|\K(\Sk)|$ at $k=75,51,31$ is a float LU determinant
(\texttt{slogdet}) of a $260$- to $869$-dimensional matrix; the relative error in the
logarithm is at the level of machine precision times the dimension, but it is not an
exact integer and is not claimed to be.

\emph{Single-strand.} As in \cite{Bio1,Bio2,Bio3,Bio4}. Section~\ref{sec:rrna} is the
measurement of what that costs here.

\emph{Open.} (a) The two exact arrangement groups each carry one large prime factor
($21\,308\,303$ at $k=150$, $3\,448\,649$ at $k=100$) that we do not interpret; as in
\cite{Bio3} that is what the sandpile group of a generic digraph looks like, and no
topological reading is claimed. (b) The $870$-vertex skeleton at $k=31$ is out of exact
reach; a modular Smith form would settle it, and nothing in principle prevents that.
(c) Whether the $10^{344}$ growth of the arrangement summand between $k=150$ and $k=31$
has a law in terms of the repeat-length distribution.

\section*{Provenance}

Unrefereed preprint. Lemma~\ref{lem:branch} is elementary and is stated because it is
what makes the computation possible, not because it is deep; everything else in this note
is computation on public data. The genome is NC\_000913.3 and the annotation is the NCBI
feature table for the same accession, both downloaded by the script and cached with it.
The expectation of a multiplicity-seven rRNA family, which Section~\ref{sec:rrna}
corrects, was the premise of the task that prompted this note, and is recorded here
because the correction is the useful part.

\begin{thebibliography}{99}

\bibitem{BCALM} R.~Chikhi, A.~Limasset and P.~Medvedev,
\emph{Compacting de Bruijn graphs from sequencing data quickly and in low memory},
Bioinformatics \textbf{32} (2016), i201--i208.

\bibitem{Blattner} F.~R.~Blattner \emph{et al.},
\emph{The complete genome sequence of Escherichia coli K-12},
Science \textbf{277} (1997), 1453--1462.

\bibitem{Bio1} R.~Chen,
\emph{Assembly ambiguity is not determined by branch statistics: critical groups of de
Bruijn graphs}, companion note (Bio~1, version~2), 2026.

\bibitem{Bio2} R.~Chen,
\emph{The sandpile group as a coordinate system on genome reconstructions}, companion
note (Bio~2), 2026.

\bibitem{Bio3} R.~Chen,
\emph{The sandpile group of a de Bruijn graph splits along its repeats}, companion note
(Bio~3), 2026.

\bibitem{Bio4} R.~Chen,
\emph{Multi-copy repeats: where the two-copy theory of assembly ambiguity stops},
companion note (Bio~4), 2026.

\bibitem{Flye} M.~Kolmogorov, J.~Yuan, Y.~Lin and P.~A.~Pevzner,
\emph{Assembly of long, error-prone reads using repeat graphs},
Nature Biotechnology \textbf{37} (2019), 540--546.

\bibitem{KM} J.~D.~Kececioglu and E.~W.~Myers,
\emph{Combinatorial algorithms for DNA sequence assembly},
Algorithmica \textbf{13} (1995), 7--51.

\bibitem{PTT} P.~A.~Pevzner, H.~Tang and G.~Tesler,
\emph{De novo repeat classification and fragment assembly},
Genome Research \textbf{14} (2004), 1786--1796.

\bibitem{PTW} P.~A.~Pevzner, H.~Tang and M.~S.~Waterman,
\emph{An Eulerian path approach to DNA fragment assembly},
Proc.\ Natl.\ Acad.\ Sci.\ USA \textbf{98} (2001), 9748--9753.

\end{thebibliography}

\end{document}
